\oldpage[IV-4]{80}
où $Y=\Spec(A)$, $A$ étant un anneau local. Comme $f$ est de présentation finie, on a $X=\Spec(B)$, où $B$ est un $A$-module plat de présentation finie (1.4.7), donc \emph{libre} (Bourbaki, \emph{Alg. comm.}, chap. II, § 5, n\textsuperscript{o} 2, cor. 2 du th. 1). En outre, si $\mathfrak m$ est l'idéal maximal de $A$ et $k$ son corps résiduel, $B/\mathfrak mB$ est par hypothèse un corps, à la fois extension radicielle et extension finie séparable de $k$, puisque $f$ est étale et radiciel (17.6.1); donc $B/\mathfrak mB$ est isomorphe à $k$, et comme $B$ est un $A$-module libre, $B$ est isomorphe à $A$. C.Q.F.D.

\begin{corollary}[17.9.2]
\label{IV.17.9.2}
Soit $X$ un préschéma connexe; si $f:X\to Y$ est une immersion fermée étale, alors $f$ est un isomorphisme de $X$ sur une composante connexe ouverte de $Y$.
\end{corollary}

En effet $f$ est étale et radiciel, donc un isomorphisme de $X$ sur un préschéma induit sur une partie ouverte de $Y$; mais par hypothèse $f(X)$ est fermé dans $Y$, donc à la fois ouvert et fermé, et puisque $f(X)$ est connexe, c'est une composante connexe de $Y$.

\begin{corollary}[17.9.3]
\label{IV.17.9.3}
Soit $f:X\to Y$ un morphisme étale \textup{(resp.}\ étale et séparé\textup{)}. Alors toute $Y$-section $g:Y\to X$ de $X$ est une immersion ouverte \textup{(resp.}\ ouverte et fermée\textup{)}; en outre l'application $g\leadsto g(Y)$ est une bijection de l'ensemble $\Gamma(X/Y)$ des $Y$-sections de $X$ sur l'ensemble des parties ouvertes $Z$ \textup{(resp.}\ ouvertes et fermées\textup{)} de $X$ telles que la restriction de $f$ à $Z$ soit un morphisme surjectif et radiciel de $Z$ sur $Y$.
\end{corollary}

En effet, le fait que $g$ soit une immersion ouverte résulte déjà de ce que $f$ est non ramifié (17.4.1, \emph{b$''$)}), et la restriction de $f$ à l'ouvert $g(Y)$ de $X$ est un isomorphisme. Inversement, si $Z$ est un ouvert de $X$ tel que $f|Z$ soit un morphisme surjectif et radiciel de $Z$ sur $Y$, $f|Z$ est un isomorphisme en vertu de (17.9.1), puisqu'il est étale. Si $f$ est en outre séparé, on sait que $g$ est une immersion fermée (\textbf{I}, 5.4.6), ce qui achève de prouver le corollaire.

\begin{corollary}[17.9.4]
\label{IV.17.9.4}
Soient $Y$ un préschéma connexe, $f:X\to Y$ un morphisme étale et séparé. Alors toute $Y$-section $g$ de $X$ est un isomorphisme de $Y$ sur une composante connexe ouverte de $X$, et l'application $g\leadsto g(Y)$ est une bijection de $\Gamma(X/Y)$ sur l'ensemble des composantes connexes ouvertes $Z$ de $X$ telles que la restriction de $f$ à $Z$ soit un morphisme surjectif et radiciel de $Z$ sur $Y$.
\end{corollary}

\begin{corollary}[17.9.5]
\label{IV.17.9.5}
Soient $g:Y\to S$, $h:X\to S$ deux morphismes localement de présentation finie; on suppose en outre que $h$ est \textup{plat}. Pour qu'un $S$-morphisme $f:X\to Y$ soit une immersion ouverte \textup{(resp.}\ un isomorphisme\textup{)}, il faut et il suffit que, pour tout $s\in S$, le morphisme $f_s:h^{-1}(s)\to g^{-1}(s)$ déduit de $f$ par le changement de base $\Spec(\boldsymbol{k}(s))\to S$, soit une immersion ouverte \textup{(resp.}\ un isomorphisme\textup{)}.
\end{corollary}

En effet, si $f_s$ est une immersion ouverte pour tout $s\in S$, il résulte de (17.8.3) que $f$ est un morphisme étale; comme pour tout $y\in Y$, on a $f^{-1}(y)=f_s^{-1}(y)$ avec $s=g(y)$, $f$ est radiciel; donc $f$ est une immersion ouverte en vertu de (17.9.1). Si de plus $f_s$ est surjectif pour tout $s\in S$, $f$ est surjectif, donc un isomorphisme.

La proposition suivante précise (10.4.11):

\begin{proposition}[17.9.6]
\EGAReaderTarget{ega4.c04.prop.17-9-6}{IV.17.9.6}
Soient $S$ un préschéma, $X$ un $S$-préschéma de présentation finie. Tout $S$-endomorphisme de $X$ qui est un monomorphisme est un automorphisme de $X$.
\end{proposition}

Soient $f:X\to S$ le morphisme structural, $g$ le $S$-endomorphisme considéré. La question est locale sur $S$, et on peut par suite supposer que $S$ est affine. Utilisant (8.9.1) et (8.10.5, (i \emph{bis})), on est ramené au cas où $S=\Spec(A)$, où $A$ est une $\mathbb{Z}$-algèbre de type fini, et par suite $X$ est un $\mathbb{Z}$-préschéma de type fini. Il résulte déjà de (10.4.11) que $g$ est un morphisme bijectif, puisqu'un monomorphisme est radiciel (8.11.5.1); il suffira donc de montrer
\oldpage[IV-4]{81}
que $g$ est une immersion ouverte, et puisque $g$ est radiciel, il suffira, en vertu de (17.9.1), de prouver que $g$ est \emph{étale}. En outre, comme l'ensemble des points où $g$ est étale est ouvert et que $X$ est un préschéma de Jacobson (10.4.7), il suffira de montrer que $g$ est étale en tout point \emph{fermé} $z$ de $X$ (10.3.1). Posons $z'=g(z)$; il résulte de (10.4.11.1, (i)) que $z'$ est aussi un point fermé de $X$ et puisque $g$ est un monomorphisme, l'application $\boldsymbol{k}(z')\to\boldsymbol{k}(z)$ déduite de $g$ est un \emph{isomorphisme}. Pour que $g$ soit étale au point $z$, il faut et il suffit donc que l'homomorphisme canonique $\widehat{\sh{O}}_{X,z'}\to\widehat{\sh{O}}_{X,z}$ soit bijectif (17.6.3, \emph{e$''$)}). Nous allons pour cela prouver que pour tout entier $n$, l'homomorphisme canonique $\sh{O}_{X,z'}/\mathfrak m_{z'}^{n+1}\to\sh{O}_{X,z}/\mathfrak m_z^{n+1}$ est \emph{bijectif}, d'où résultera aussitôt la conclusion. On peut supposer que $z$ appartient à l'ensemble fini $T_{p,d}$ des points fermés $t$ de $X$ tels que $\boldsymbol{k}(t)$ soit une extension de $\mathbb{F}_p$ dont le degré divise $d$; auquel cas il en est de même de $z'$. Désignons par $\sh{J}$ l'Idéal cohérent de $\sh{O}_X$ tel que $\sh{J}_x=\sh{O}_x$ pour $x\notin T_{p,d}$, et $\sh{J}_x=\mathfrak m_x^{n+1}$ pour $x\in T_{p,d}$. Soient $T_{p,d,n}$ le sous-préschéma fermé de $X$ défini par $\sh{J}$, $j:T_{p,d,n}\to X$ l'injection canonique; le morphisme composé $g\circ j:T_{p,d,n}\to X$ applique l'ensemble $T_{p,d}$ dans lui-même, et pour tout $x\in T_{p,d}$, l'homomorphisme $\sh{O}_{X,g(x)}\to\sh{O}_{X,x}/\mathfrak m^{n+1}$ se factorise en
\[
\sh{O}_{X,g(x)}\to\sh{O}_{X,g(x)}/\mathfrak m_{g(x)}^{n+1}\to\sh{O}_{X,x}/\mathfrak m_x^{n+1}.
\]
Donc (\textbf{I}, 4.1.9) il existe un unique endomorphisme $g'$ de $T_{p,d,n}$ tel que $j\circ g'=g\circ j$. Comme $j$ est un monomorphisme, et qu'il en est de même de $g$ par hypothèse, on en déduit que $g'$ est aussi un monomorphisme, et la proposition sera prouvée si l'on montre que $g'$ est un \emph{automorphisme} de $T_{p,d,n}$. Or, pour tout $x\in T_{p,d,n}$, on a vu que $\boldsymbol{k}(x)$ est un corps fini, et comme $\sh{O}_{X,x}$ est noethérien, chacun des $\mathfrak m_x^h/\mathfrak m_x^{h+1}$ est un $\boldsymbol{k}(x)$-espace vectoriel de rang fini; d'où l'on conclut aussitôt que l'anneau local $\sh{O}_{X,x}/\mathfrak m_x^{n+1}$ de $T_{p,d,n}$ au point $x$ a un nombre \emph{fini} d'éléments, et \emph{a fortiori} est un $\mathbb{Z}$-module de type fini. Comme $T_{p,d,n}$ est somme des préschémas $\Spec(\sh{O}_{X,x}/\mathfrak m_x^{n+1})$ en nombre fini, c'est un $\mathbb{Z}$-préschéma \emph{fini}; l'endomorphisme $g'$ est donc lui aussi \emph{fini} (\textbf{II}, 6.1.\textbf{5}, (v)), donc propre. Mais un monomorphisme propre est une \emph{immersion fermée} (8.11.5); si $T_{p,d,n}=\Spec(B)$, $g'$ correspond donc à un endomorphisme surjectif $\varphi:B\to B$ de l'anneau $B$. Or, l'ensemble $B$ est fini, donc $\varphi$ est nécessairement bijectif. C.Q.F.D.

\subsection{Dimension relative d'un préschéma lisse sur un autre.}
\EGAReaderTarget{ega4.c04.subsec.17-10}{subsection:IV.17.10}

\begin{definition}[17.10.1]
\EGAReaderTarget{ega4.c04.def.17-10-1}{IV.17.10.1}
Soit $f:X\to Y$ un morphisme localement de type fini. On appelle dimension relative de $f$ au point $x\in X$ (ou dimension relative de $X$ sur $Y$ au point $x$) et l'on note $\dim_x f$ l'entier positif $\dim_x(f^{-1}(f(x)))$.
\end{definition}

Dire que $f$ est \emph{quasi-fini} au point $x$ (\textbf{Err}$_{\mathrm{III}}$, 20) équivaut donc à dire que $\dim_x f=0$. On a vu (13.1.3) que la fonction $x\leadsto\dim_x f$ est semi-continue supérieurement. On notera que, même lorsque le morphisme $f$ a la propriété $(S_1)$ (autrement dit (6.8.1) est plat et tel que ses fibres n'aient pas de cycle premier associé immergé), la fonction $x\leadsto\dim_x f$ n'est pas nécessairement continue, comme le montre l'exemple où $Y=\Spec(k)$, où $k$ est un corps, et $X=\Spec(k[U,V,W]/\mathfrak p\mathfrak q)$, où $\mathfrak p=(W)$ et $\mathfrak q=(U)+(V\setmin W)$, idéaux premiers de $k[U,V,W]$ ($X$ étant donc la réunion, dans l'espace à 3 dimensions, d'un plan et d'une droite non parallèle à ce plan).

Dire qu'un morphisme localement de présentation finie $f:X\to Y$ est étale au point $x\in X$ signifie encore que $f$ est lisse au point $x$ et que $\dim_x f=0$ (17.6.1).

\begin{proposition}[17.10.2]
\label{IV.17.10.2}
Soit $f:X\to Y$ un morphisme lisse. Pour tout $x\in X$, le $\sh{O}_{X,x}$-Module localement libre $\varOmega_f^1$ \textup{(17.2.3)} a un rang en $x$ égal à $\dim_x f$ \textup{(ce qui entraîne que $x\leadsto\dim_x f$ est une fonction \emph{continue} dans $X$)}.
\end{proposition}

En effet, si $y=f(x)$, $X_y=f^{-1}(y)$ est lisse sur $\boldsymbol{k}(y)$, et si $f_y:X_y\to\Spec(\boldsymbol{k}(y))$ est le morphisme structural, on a $\varOmega_{f_y}^1=\varOmega_f^1\otimes_{\sh{O}_{Y,y}}\boldsymbol{k}(y)$ (16.4.5); on peut donc se borner au cas où $Y=\Spec(k)$ est le spectre d'un corps; en outre, en vertu de (16.4.5) et (17.7.1), on peut remplacer $k$ par une extension algébriquement close, autrement dit
\oldpage[IV-4]{82}
supposer $k$ \emph{algébriquement clos}. L'ensemble des points de $X$ rationnels sur $k$ étant alors dense dans $X$ (10.4.8), on peut se borner au cas où $x$ est rationnel sur $k$. Or (16.4.12) $(\varOmega^1_{X/k})_x\otimes_{\sh{O}_x}\boldsymbol{k}(x)$ est alors $k$-isomorphe à $\mathfrak m_x/\mathfrak m_x^2$, et puisque $\sh{O}_x$ est un anneau local régulier (17.5.1), $\operatorname{rg}_k(\mathfrak m_x/\mathfrak m_x^2)=\dim(\sh{O}_x)$ (\textbf{0}, 17.1.1); donc $\operatorname{rg}_{\sh{O}_x}(\varOmega^1_{X/k})_x=\dim(\sh{O}_x)$. Mais puisque $\boldsymbol{k}(x)=k$, on a $\dim_x f=\dim(\sh{O}_x)$ (5.2.3). C.Q.F.D.

Nous démontrerons plus loin une réciproque de ce résultat (17.15.5).

\begin{corollary}[17.10.3]
\label{IV.17.10.3}
Soient $f:X\to Y$, $g:Y\to Z$ deux morphismes lisses. Alors, pour tout $x\in X$, on a
\begin{equation}
\label{IV.17.10.3.1}
\tag{17.10.3.1}
\dim_x(g\circ f)=\dim_x f+\dim_{f(x)}g.
\end{equation}
\end{corollary}

En effet, $g\circ f$ est lisse (17.3.3), donc les trois $\sh{O}_X$-Modules $\varOmega^1_{X/Y}$, $\varOmega^1_{X/Z}$ et $f^*(\varOmega^1_{Y/Z})$ sont localement libres (17.2.3 et $\mathbf{0}_{\mathrm I}$, 5.4.5); en outre le rang en $x$ de $f^*(\varOmega^1_{Y/Z})$ est égal au rang en $y=f(x)$ de $\varOmega^1_{Y/Z}$. L'égalité (17.10.3.1) est donc conséquence de (17.10.2) et de l'exactitude de la suite (17.2.3.1).

\begin{corollary}[17.10.4]
\label{IV.17.10.4}
Soient $f:X\to Y$ un morphisme lisse, $X'$ un sous-préschéma de $X$ tel que le morphisme composé $X'\xrightarrow{j}X\xrightarrow{f}Y$ (où $j$ est l'injection canonique) soit lisse. Alors le faisceau conormal $\sh{N}_{X'/X}$ est un $\sh{O}_{X'}$-Module localement libre, et pour tout $x\in X'$, on a
\begin{equation}
\label{IV.17.10.4.1}
\tag{17.10.4.1}
\dim_x f=\dim_x(f\circ j)+\operatorname{rg}_{\sh{O}_{X',x}}(\sh{N}_{X'/X})_x.
\end{equation}
\end{corollary}

En effet, $\varOmega^1_{X/Y}\otimes_{\sh{O}_X}\sh{O}_{X'}$ et $\varOmega^1_{X'/Y}$ sont alors localement libres et la suite exacte (17.2.5.1) est scindée dans un voisinage convenable de chaque point de $X'$, donc $\sh{N}_{X'/X}$ est localement libre (Bourbaki, \emph{Alg. comm.}, chap. II, § 5, n\textsuperscript{o} 2, th. 1), et la relation (17.10.4.1) résulte aussitôt de l'exactitude de la suite (17.2.5.1).

\subsection{Morphismes lisses de préschémas lisses.}
\EGAReaderTarget{ega4.c04.subsec.17-11}{subsection:IV.17.11}

\begin{theorem}[17.11.1]
\label{IV.17.11.1}
Soient $f:Y\to S$ et $h:X\to S$ deux morphismes localement de présentation finie, $g:X\to Y$ un $S$-morphisme, $x$ un point de $X$; posons $y=g(x)$, $s=f(y)=h(x)$. Les conditions suivantes sont équivalentes:
\begin{enumerate}
  \item[\textup{a)}] $f$ est lisse au point $y$ et $g$ est lisse au point $x$.
  \item[\textup{b)}] $g$ et $h$ sont lisses au point $x$.
  \item[\textup{c)}] $h$ est lisse au point $x$, et l'homomorphisme canonique \textup{(16.4.18)}
  \begin{equation}
  \label{IV.17.11.1.1}
  \tag{17.11.1.1}
  (g^*(\varOmega^1_{Y/S}))_x\to(\varOmega^1_{X/S})_x
  \end{equation}
  est inversible à gauche \textup{(autrement dit est un isomorphisme sur un \emph{facteur direct de} $(\varOmega^1_{X/S})_x$)}.
  \item[\textup{c$'$)}] $h$ est lisse au point $x$, et l'homomorphisme canonique
  \begin{equation}
  \label{IV.17.11.1.2}
  \tag{17.11.1.2}
  ((\varOmega^1_{Y/S})_y\otimes_{\sh{O}_{Y,y}}\boldsymbol{k}(y))\otimes_{\mathbf{k}(y)}\boldsymbol{k}(x)\to(\varOmega^1_{X/S})_x\otimes_{\sh{O}_{X,x}}\boldsymbol{k}(x)
  \end{equation}
  est injectif.
\end{enumerate}

Supposons en outre que l'homomorphisme $\boldsymbol{k}(y)\to\boldsymbol{k}(x)$ soit bijectif. Alors les conditions précédentes sont aussi équivalentes à la suivante:
\oldpage[IV-4]{83}
\begin{enumerate}
  \item[\textup{d)}] $h$ est lisse au point $x$, et l'application canonique $\mathrm{T}_{X/S}(x)\to\mathrm{T}_{Y/S}(y)$ de l'espace vectoriel tangent en $x$ à $X$ dans l'espace vectoriel tangent en $y$ à $Y$ \textup{(16.5.12)} est surjective.
\end{enumerate}
\end{theorem}

Le fait que \emph{a)} entraîne \emph{b)} résulte trivialement de (17.3.3, (ii)); \emph{b)} entraîne \emph{c)} par application de (17.2.3, (ii)). Pour voir que \emph{c)} est équivalente à \emph{c')}, notons qu'en vertu de (17.2.3, (i)), $\varOmega^1_{X/S}$ est un $\sh{O}_X$-Module localement libre de type fini; il suffit alors d'appliquer (\textbf{0}, 19.1.12) à l'anneau local $\sh{O}_{X,x}$ et à l'homomorphisme (17.11.1.1) de $\sh{O}_{X,x}$-modules de type fini. Lorsque $\boldsymbol{k}(x)=\boldsymbol{k}(y)$, l'application linéaire tangente $\mathrm{T}_{X/S}(x)\to\mathrm{T}_{Y/S}(y)$ est transposée de (17.11.1.2), par (16.5.12), ce qui établit dans ce cas l'équivalence de \emph{c')} et \emph{d)}.

Il reste donc à prouver que \emph{c)} entraîne \emph{a)}. On peut se borner au cas où $S=\Spec(A)$, $Y=\Spec(B)$, $X=\Spec(C)$ sont affines. L'hypothèse implique (17.2.3, (i)) que $(\varOmega^1_{X/S})_x=\Omega^1_{C_x/A_s}$ est un $C_x$-module libre: en effet (16.10.6) il existe des éléments $t_i$ ($1\leq i\leq r$) de $B_y=\sh{O}_{Y,y}$ tels que les différentielles $d_{B_y/A_s}(t_i)$ engendrent le $B_y$-module $\Omega^1_{B_y/A_s}$ et que leurs images dans $\Omega^1_{C_x/A_s}$ fassent partie d'une base de ce $C_x$-module libre. Comme $\varOmega^1_{Y/S}$ (resp. $\varOmega^1_{X/S}$) est un $\sh{O}_Y$-Module (resp. un $\sh{O}_X$-Module) de présentation finie (16.4.22), on peut, en remplaçant au besoin $X$ et $Y$ par des voisinages ouverts affines convenables de $x$ et $y$ respectivement, supposer que $\Omega^1_{C/A}$ est un $C$-module libre et que les $t_i$ sont les images d'éléments $s_i$ ($1\leq i\leq r$) de $B$ tels que les $d_{B/A}(s_i)$ engendrent le $B$-module $\Omega^1_{B/A}$ et que leurs images dans $\Omega^1_{C/A}$ fassent partie d'une base de ce $C$-module (Bourbaki, \emph{Alg. comm.}, chap. II, § 5, n\textsuperscript{o} 1, prop. 2). Soit $\varphi$ le $A$-homomorphisme de $B'=A[T_1,\ldots,T_r]$ dans $B$ tel que $\varphi(T_i)=s_i$ pour tout $i$; le di-homomorphisme correspondant $\Omega^1_{B'/A}\to\Omega^1_{B/A}$ (\textbf{0}, 20.5.2) transforme les $d_{B'/A}(T_i)$, qui forment une base de $\Omega^1_{B'/A}$ (\textbf{0}, 20.4.13), en les $d_{B/A}(s_i)$ et est par suite \emph{surjectif}; si $Y'=\Spec(B')$, et si $u:Y\to Y'$ est le $S$-morphisme correspondant à $\varphi$, on conclut de (17.2.2) que $u$ est \emph{non ramifié}. Si l'on prouve que le morphisme composé $u\circ g:X\to Y'$ est lisse au point $x$, il résultera donc de (17.7.10) que $u$ est \emph{étale} au point $y$, puis de (17.3.5) que $g$ est lisse au point $x$; enfin, comme le morphisme structural $f':Y'\to S$ est lisse (17.3.8), $f=f'\circ u$ sera lisse au point $y$. Comme les images canoniques des $d_{B'/A}(T_i)$ dans $\Omega^1_{C/A}$ sont celles des $d_{B/A}(s_i)$, elles font partie d'une base de $\Omega^1_{C/A}$. On voit ainsi que pour prouver que \emph{c)} implique \emph{a)}, on peut se borner au cas où $Y'=Y$, et par suite supposer que $f$ est un morphisme \emph{lisse}.

En vertu de (17.8.2), on peut alors se borner au cas où $S=\Spec(k)$ est le spectre d'un corps; de plus, grâce à (17.7.1, (ii)), on peut supposer que $k$ est algébriquement clos; enfin, remplaçant au besoin $X$ et $Y$ par des voisinages ouverts de $x$ et $y$ respectivement, on peut supposer que $f$ et $h$ sont tous deux lisses, et que l'homomorphisme canonique
\[
g^*(\varOmega^1_{Y/S})\to\varOmega^1_{X/S}
\]
est inversible à gauche (\textbf{0}, 19.1.12). Alors l'ensemble des points de $X$ rationnels sur $k$ est très dense dans $X$ (10.4.8), donc, pour prouver que $g$ est lisse, il suffit de prouver que $g$ est lisse en tout point $x\in X$ \emph{rationnel sur $k$}, ou encore qu'en un tel point, $\sh{O}_{X,x}$ est un $\sh{O}_{Y,y}$-module plat et $\sh{O}_{X,x}/\mathfrak m_y\sh{O}_{X,x}$ un anneau régulier (17.5.1). Or, $y=g(x)$ est
\oldpage[IV-4]{84}
\emph{a fortiori} rationnel sur $k$, et comme $\sh{O}_{Y,y}$ est alors une $k$-algèbre formellement lisse pour les topologies discrètes et que $\sh{O}_{Y,y}/\mathfrak m_y=k$, il résulte de (\textbf{0}, 20.5.14) que l'homomorphisme canonique $\mathfrak m_y/\mathfrak m_y^2\to(\varOmega^1_{Y/S})_y\otimes_{\sh{O}_{Y,y}}\boldsymbol{k}(y)$ est bijectif; de même, l'homomorphisme canonique $\mathfrak m_x/\mathfrak m_x^2\to(\varOmega^1_{X/S})_x\otimes_{\sh{O}_{Y,x}}\boldsymbol{k}(x)$ est bijectif. L'hypothèse \emph{c')}, équivalente à \emph{c)}, signifie donc ici que l'homomorphisme canonique
\[
(\mathfrak m_y/\mathfrak m_y^2)\otimes_{\mathbf{k}(y)}\boldsymbol{k}(x)\to\mathfrak m_x/\mathfrak m_x^2
\]
est injectif. Comme l'anneau $\sh{O}_{X,x}$ est régulier, la conclusion résulte de (\textbf{0}, 17.3.3). C.Q.F.D.

\begin{corollary}[17.11.2]
\label{IV.17.11.2}
Sous les hypothèses générales de \textup{(17.11.1)}, les conditions suivantes sont équivalentes:
\begin{enumerate}
  \item[\textup{a)}] $f$ est lisse au point $y$ et $g$ est étale au point $x$.
  \item[\textup{b)}] $h$ est lisse au point $x$ et $g$ est étale au point $x$.
  \item[\textup{c)}] $h$ est lisse au point $x$, et l'homomorphisme canonique
  \[
  (g^*(\varOmega^1_{X/S}))_x\to(\varOmega^1_{X/S})_x
  \]
  est bijectif.
  \item[\textup{c$'$)}] $h$ est lisse au point $x$, et l'homomorphisme canonique
  \[
  ((\varOmega^1_{Y/S})_y\otimes_{\sh{O}_{Y,y}}\boldsymbol{k}(y))\otimes_{\mathbf{k}(y)}\boldsymbol{k}(x)\to(\varOmega^1_{X/S})_x\otimes_{\sh{O}_{X,x}}\boldsymbol{k}(x)
  \]
  est bijectif.
\end{enumerate}

Supposons en outre que l'homomorphisme $\boldsymbol{k}(y)\to\boldsymbol{k}(x)$ soit bijectif. Alors les conditions précédentes sont aussi équivalentes à la suivante:
\begin{enumerate}
  \item[\textup{d)}] $h$ est lisse au point $x$, et l'application canonique $\mathrm{T}_{X/S}(x)\to\mathrm{T}_{Y/S}(y)$ \textup{(16.5.12)} est bijective.
\end{enumerate}
\end{corollary}

Chacune des conditions \emph{a)}, \emph{b)}, \emph{c)} de (17.11.2) est équivalente à la conjonction de la condition correspondante de (17.11.1) et du fait que $g$ est non ramifié au point $x$, compte tenu de (17.2.2) en ce qui concerne la condition \emph{c)}; d'où l'équivalence de \emph{a)}, \emph{b)} et \emph{c)}. L'équivalence de \emph{c)} et \emph{c')} résulte de l'équivalence des conditions correspondantes de (17.11.1) et du lemme de Nakayama; l'équivalence de \emph{c')} et \emph{d)} lorsque $\boldsymbol{k}(x)=\boldsymbol{k}(y)$ est immédiate par transposition (16.5.12).

\begin{corollary}[17.11.3]
\label{IV.17.11.3}
Soient $h:X\to S$ un morphisme lisse, $s_i$ \textup{($1\leq i\leq r$)} des sections de $\sh{O}_X$ au-dessus de $X$ \textup{(}qui sont aussi des sections de $h_*(\sh{O}_X)$ au-dessus de $S$\textup{)}, $g:X\to S[T_1,\ldots,T_r]=\mathbf{V}_S^r$ le $S$-morphisme correspondant à l'homomorphisme de $\sh{O}_S$-Modules $\sh{O}_S^r\to h_*(\sh{O}_X)$ défini par ces sections \textup{(\textbf{II}, 1.2.7)}. Pour que $g$ soit lisse \textup{(resp.}\ étale\textup{)} en un point $x\in X$ il faut et il suffit que les $(d_{X/S}(s_i))_x$ fassent partie d'une base \textup{(resp.}\ constituent une base\textup{)} du $\sh{O}_{X,x}$-module $(\varOmega^1_{X/S})_x$.
\end{corollary}

Il suffit d'appliquer (17.11.1) (resp. (17.11.2)) en prenant pour $f$ le morphisme structural $S[T_1,\ldots,T_r]\to S$.

\begin{corollary}[17.11.4]
\label{IV.17.11.4}
Pour qu'un morphisme $h:X\to S$ soit lisse en un point $x\in X$, il faut et il suffit qu'il existe un voisinage ouvert $U$ de $x$, un entier $r$, et un $S$-morphisme étale $U\to S[T_1,\ldots,T_n]$.
\end{corollary}

